\answer
\begin{enumerate}

%%--- 章末問題 2.1 ---
\item[\ref{ex2.1}]
$\exp(ev_D/kT) \gg 1$のとき$i_D \approx I_0\exp(ev_D/kT)$。
電流が10倍になる電圧増分$\mathit{\Delta}v$は
\begin{equation*}
\exp\left(\frac{e\mathit{\Delta}v}{kT}\right) = 10
\quad\Rightarrow\quad
\mathit{\Delta}v = \frac{kT}{e}\ln 10 = 0.026\times2.30 \approx 0.060\,\mathrm{V}
\end{equation*}
すなわちわずか60\,mVで電流は1けた変わる。
順方向電圧が0.6〜0.7\,V付近でほぼ一定に見えるのは，
この鋭い指数特性のためである。

%%--- 章末問題 2.2 ---
\item[\ref{ex2.2}]
コレクタに到達する電子は$99.5$\,\%だから
$i_C = 0.995\times40 = 39.8$\,A。
ベース電流は再結合分で$i_B = 0.005\times40 = 0.2$\,A。
式\ref{eq2.4}の$i_E = i_C + i_B = 39.8 + 0.2 = 40$\,Aも確かに成り立つ。
電流増幅率は式\ref{eq2.5}より
\begin{equation*}
\beta = \frac{i_C}{i_B} = \frac{39.8}{0.2} = 199
\end{equation*}

%%--- 章末問題 2.3 ---
\item[\ref{ex2.3}]
(a) $R_{on}{I^*}^2 = V_{on}I^*$より
\begin{equation*}
I^* = \frac{V_{on}}{R_{on}} = \frac{1.6}{0.08} = 20\,\mathrm{A}
\end{equation*}
(b) MOSFETの損失は$I^2$に，IGBTの損失は$I$に比例するから，
$I < 20$\,AではMOSFETが，$I > 20$\,AではIGBTが低損失である。
小電流・高周波はMOSFET，大電流はIGBTという使い分け（\ref{fig2.1}）は，
この損失の質の違いから生まれる。

%%--- 章末問題 2.4 ---
\item[\ref{ex2.4}]
MOSFETのオン抵抗の主因は，耐圧のために低濃度・厚めに作られた
ドリフト層の抵抗である。
IGBTではオン時に裏面のp形コレクタ層からドリフト層へ正孔が注入され，
電気的中性を保つように電子も増えるため，
層内のキャリア濃度$n$，$p$がドーピング濃度よりけた違いに大きくなる。
式\ref{eq1.5}より導電率$\sigma = e(n\mu_n + p\mu_p)$はキャリア濃度に
比例するから，ドリフト層の抵抗は大幅に下がる。これが伝導度変調である。

代償はターンオフに現れる。
ドリフト層にたまった大量の少数キャリアが再結合して消えるまで
電流が流れ続け（テール電流），
そのあいだ素子には高い電圧と電流が同時にかかるため
スイッチング損失が大きくなる。

%%--- 章末問題 2.5 ---
\item[\ref{ex2.5}]
IGBTのゲートは酸化膜の上にあり，電場でp形表面に反転層を作って
「電子の供給路」を開け閉めする。
ゲート電圧を切れば反転層が消えて電子の供給が絶たれ，
それに伴い裏面からの正孔注入も止まるので，素子は自らオフできる
（内部のnpn部分が単独で動き続けないよう，
寄生サイリスタがラッチしない設計になっている）。

一方サイリスタのゲートは内側のp層に直接つながっており，
ゲート電流で2つのトランジスタの正帰還ループに火をつける。
いったんループが自立すると，
主電流のキャリア注入そのものが駆動電流を賄うため，
ゲートからはもはや介入できない。
「電場で通り道を作る」IGBTと「電流で正帰還を点火する」サイリスタ ―
同じ4層でも，制御のしくみの違いがオフの可否を分けている。

%%--- 章末問題 2.6 ---
\item[\ref{ex2.6}]
\rev{（シミュレーションの確認ポイント）
(a) 順方向は$0.6$〜$0.7$\,V付近から電流が急に立ち上がり，
逆方向は逆方向飽和電流$I_0$程度（\texttt{1N4148}のモデルではnA規模）しか流れず，
リニア軸ではほぼゼロに見える。式\ref{eq2.2}の指数特性そのものである。}

\rev{(b) 縦軸を対数にすると順方向の特性は直線になり，その傾きが「電流を10倍にするのに要する電圧」である。
式\ref{eq2.2}のとおりの理想的なダイオード（モデルを指定しない既定の\texttt{diode}）では，
章末問題\ref{ex2.1}の答えと同じ1けたあたり約$60$\,mVになる。
\texttt{1N4148}などの実素子のモデルでは，指数が$\exp(ev_D/nkT)$の形に
理想係数$n$（$1$〜$2$程度）を含むため，傾きは$n\times60$\,mV
（\texttt{1N4148}のモデルでは$n \approx 1.75$で約$100$\,mV）と読める。
また大電流側では直列抵抗の電圧降下のぶん直線から外れ，電流の伸びが鈍る。}

\rev{(c) 出力特性は原点からまっすぐ立ち上がり（ダイオードのような$0.7$\,Vの下駄がない），
原点付近の直線の傾きの逆数$v_{DS}/i_D$がオン抵抗$R_{on}$である（例題\ref{rei2.2a}）。
$v_{GS}$を$4$，$6$，$8$，$10$\,Vと上げるほど傾きは急になり，$R_{on}$は小さくなる。
数値はモデルによるが，\texttt{IRF540}のモデルではしきい値電圧が約$4$\,Vなので
$v_{GS} = 4$\,Vではほとんど電流が流れず，$6$\,Vで$0.1\,\Omega$程度，
$10$\,Vで$0.04\,\Omega$程度まで下がる。
スイッチとして使うときにゲートをしきい値よりずっと高く駆動する理由が，この図に表れている。}

\end{enumerate}
